% part: intuitionistic-logic
% chapter: semantics
% section: introduction

\documentclass[../../../include/open-logic-section]{subfiles}

\begin{document}

\olfileid{int}{sem}{int}

\olsection{مقدمة}

لا يستوفى وصف منطق من غير دلالة، والمنطق الحدسي من ذلك. وللمنطق الكلاسيكي
دلالة معيارية قائمة على ال!!{valuation}s؛ وأما الحدسي فتتنازع تفسيره
دلالات عدة، ولا تسلم واحدة منها من الاعتراض إذا طلبنا منها تفسيرًا لمعاني
الروابط يرضاه الحدسي من كل وجه.

وأشهرها فيما يظهر الدلالة القائمة على ال!!{relational model}s، وهي شبيهة
بدلالة المنطق الجهوي. تسند فيها ال!!{propositional variable}s إلى عوالم،
وتصل العوالم علاقة نفاذ تكون دائمًا ترتيبًا جزئيًا: انعكاسية ومضادة للتناظر
ومتعدية.

وتتصور هذه العوالم، على سبيل الحدس، أحوالًا للمعرفة أو «مواقف استدلالية».
فتتاح الحالة~${\OLId{latin-ordinary}{w}}'$ من~${\OLId{latin-ordinary}{w}}$ إذا وفقط إذا كانت، بحسب ما نعلم، حالًا معرفية
مستقبلية ممكنة توافق المعلوم عند~${\OLId{latin-ordinary}{w}}$. وما علم لا يصير بعد ذلك غير معلوم؛
فإذا كانت~$!A$ معلومة عند~${\OLId{latin-ordinary}{w}}$ وكان~$Rww'$، بقيت~$!A$ معلومة عند~${\OLId{latin-ordinary}{w}}'$.
فــ«المعرفة» بهذا المعنى رتيبة بالإضافة إلى علاقة النفاذ.

ولنحد «كون~$!A$ معلومة»، كما في منطق المعرفة، بصدقها في جميع البدائل
المعرفية. فعندئذ تكون~$!A \land !B$ معلومة عند~${\OLId{latin-ordinary}{w}}$ متى كانت~$!A$ و$!B$
معلومتين معًا في جميع تلك البدائل. ولما كانت المعرفة رتيبة وكانت~${\OLId{latin-looped}{R}}$
انعكاسية، رجع ذلك إلى كون~$!A \land !B$ معلومة عند~${\OLId{latin-ordinary}{w}}$ إذا وفقط إذا كانت
$!A$ و$!B$ معلومتين عند~${\OLId{latin-ordinary}{w}}$. وعلى القياس نفسه تكون~$!A \lor !B$ معلومة
عند~${\OLId{latin-ordinary}{w}}$ إذا وفقط إذا كانت إحداهما على الأقل معلومة. فشرطا صدق~$\land$
و$\lor$ هنا هما شرطاهما في المنطق الكلاسيكي.

وأما الشرطية فيختلف شرط صدقها: تكون~$!A \lif !B$ معلومة عند~${\OLId{latin-ordinary}{w}}$ إذا وفقط
إذا لم توجد~${\OLId{latin-ordinary}{w}}'$ بحيث~$Rww'$ تعلم عندها~$!A$ ولا تعلم~$!B$. وليس هذا
بمجرد أن تكون~$!A$ غير معلومة أو~$!B$ معلومة عند~${\OLId{latin-ordinary}{w}}$؛ إذ قد نجهل~$!A$
و$!B$ عند~${\OLId{latin-ordinary}{w}}$، ثم توجد حال معرفية مستقبلية~${\OLId{latin-ordinary}{w}}'$ بحيث~$Rww'$ تعلم فيها
$!A$ ولا تعلم~$!B$.

ولا نعلم~$\lnot !A$ إلا إذا امتنع وجود حال معرفية مستقبلية ممكنة نعلم فيها
$!A$. وذلك أن إمكان علم~$!A$ يجعل~$!A$ معلومة في حال مستقبلية ممكنة؛ وحيث
لا يمكن أن نعلم~$\lfalse$، نكون في تلك الحال عالمين بــ$!A$ غير عالمين
بــ$\lfalse$.

وعلى هذا يبطل مبدأ الثالث المرفوع: فقد تكون~$!A$ غير معلومة الآن مع إمكان
علمها بعد. فإذا كانت !!a{formula}~$!A$ كذلك، لم تكن~$!A$ ولا~$\lnot !A$
معلومة، فلم تكن~$!A \lor \lnot !A$ معلومة. ومع ذلك نعلم، مثلًا،
$\lnot(!A \land \lnot !A)$؛ لأنه لا توجد حال مستقبلية نعلم فيها~$!A$
و$\lnot !A$ معًا، وهذا معلوم سواء علمنا~$!A$ أو لم نعلمها، وسواء علمنا
$\lnot !A$ أو لم نعلمها.

وليست ال!!^{relational model}s هي الدلالة الوحيدة للمنطق الحدسي؛ فثمة
دلالة طوبولوجية تفسر القضايا بمجموعات مفتوحة في فضاء طوبولوجي، وتفسر
الروابط بعمليات على تلك المجموعات؛ فمثلًا يقابل~$\land$ التقاطع.

\end{document}
